
% Kapitola 6: Právne a etické aspekty
% Stav: finálna verzia

\section{Právne a etické aspekty}

Právne a etické aspekty navrhovaného riešenia nemožno oddeliť od jeho technického návrhu. Aj keď aplikácia pracuje prevažne s bibliografickými metadátami a nie s plnými textami alebo s osobitnými kategóriami osobných údajov, stále ide o údaje, ktoré môžu byť viazané na identifikované alebo identifikovateľné fyzické osoby. Zároveň sa pri importe a ďalšom spracovaní pracuje s metadátami pochádzajúcimi z viacerých externých zdrojov, pri ktorých sa líšia podmienky opätovného použitia. Táto kapitola preto stručne vymedzuje dve oblasti: spracúvanie údajov o autoroch podľa GDPR a licenčný režim bibliografických metadát \cite{GDPR2016,CrossrefMetadataRetrieval}.

\subsection{Spracúvanie údajov o autoroch a GDPR}

Nariadenie (EÚ) 2016/679 definuje osobné údaje ako akékoľvek informácie týkajúce sa identifikovanej alebo identifikovateľnej fyzickej osoby a spracúvanie ako akúkoľvek operáciu alebo súbor operácií vykonávaných s osobnými údajmi, či už automatizovanými alebo neautomatizovanými prostriedkami \cite{GDPR2016}. V prostredí bibliografických záznamov to znamená, že meno autora, jeho identifikátor ORCID, afiliácia alebo kombinácia viacerých údajov umožňujúcich jednoznačné priradenie ku konkrétnej osobe môžu mať povahu osobných údajov. Samotná práca s publikačnými metadátami preto nestojí mimo rámca ochrany osobných údajov len preto, že ide o údaje vedeckej komunikácie.

Pre navrhovanú aplikáciu je relevantné najmä to, že GDPR nestanovuje len definície, ale aj základné zásady spracúvania osobných údajov. Medzi ne patrí zákonnosť, korektnosť a transparentnosť, účelové viazanie, minimalizácia údajov, presnosť a primeraná bezpečnosť spracovania \cite{GDPR2016}. Tieto zásady dopĺňa aj výklad EDPB k článku 25, ktorý zdôrazňuje ochranu údajov už pri návrhu systému a v predvolených nastaveniach \cite{EDPBGuidelinesDPbDD2019}. V praktickom preklade pre túto prácu to znamená, že aplikácia by mala spracúvať iba tie údaje o autoroch, ktoré sú potrebné pre evidenciu a normalizáciu publikačných záznamov, mala by umožniť opravu nepresných údajov a mala by podporovať kontrolovaný prístup k rozhraniu a auditovateľnosť zásahov. Keďže predmetom spracovania nie sú súkromné texty autorov, ale najmä verejne dostupné bibliografické údaje, jadrom problému nie je utajenie obsahu, ale korektné a primerané nakladanie s identifikačnými údajmi a s históriou ich úprav.

Z etického hľadiska je dôležité aj to, aby automatizovaná podpora normalizácie neviedla k neodôvodnenému prepájaniu alebo zámene autorov. Chybná identifikácia interného autora, nesprávne priradenie afiliácie alebo nekritické prevzatie návrhu modelu môže mať dopad na evidenciu publikačnej činnosti a na jej ďalšie využitie v interných procesoch univerzity. Preto má význam, aby aj v prípadoch, keď sú vstupné údaje verejne dostupné, zostal zachovaný princíp ľudského schvaľovania a spätnej dohľadateľnosti rozhodnutí.

\subsection{Licencovanie a použiteľnosť bibliografických metadát}

Pri práci s externými zdrojmi je potrebné odlíšiť technickú dostupnosť metadát od právnych podmienok ich ďalšieho použitia. V prípade Crossref je situácia pomerne priaznivá. Oficiálna dokumentácia uvádza, že metadáta sprístupňované cez Crossref sú voľne dostupné komunite, že väčšina bibliografických metadát je považovaná za fakty nepodliehajúce autorskému právu a že Crossref-generované údaje sprístupňuje ako public domain material, pričom špecifická pozornosť sa venuje najmä abstraktom, ktoré môžu zostať predmetom autorskoprávnej ochrany vydavateľa alebo autora \cite{CrossrefMetadataRetrieval}. Pre navrhovanú aplikáciu z toho vyplýva, že použitie Crossref ako doplnkového zdroja na overovanie a dopĺňanie bibliografických polí je z licenčného hľadiska podstatne menej problematické než pri uzavretých komerčných databázach.

Odlišná situácia platí pri platformách Web of Science a Scopus. Hoci tieto služby poskytujú používateľom možnosť exportovať bibliografické záznamy, podmienky ich ďalšieho použitia sa odvíjajú od licenčných a zmluvných pravidiel konkrétnej služby a inštitucionálneho prístupu. Clarivate viaže používanie svojich produktov a služieb na všeobecné podmienky a ďalšie produktovo špecifické podmienky \cite{ClarivateTermsOfUse}. Elsevier pri vývojárskom prístupe výslovne uvádza, že jednotlivé use cases a politiky určujú, ako možno dáta získané cez jeho API používať, a že ďalšie oprávnenia závisia od typu používateľa a predplatného príslušnej služby \cite{ElsevierAPIPolicy}. Z pohľadu tejto práce je preto vhodné vychádzať z konzervatívneho predpokladu, že exportované dáta z týchto zdrojov sa majú používať v rozsahu zodpovedajúcom oprávnenému internému workflow univerzity a nie ako voľne redistribuovateľný otvorený dataset.

Rozlíšenie medzi otvorenými a licenčne viazanými zdrojmi má aj etický rozmer. Ak má aplikácia prepájať viacero zdrojov a dopĺňať chýbajúce hodnoty, musí byť zrejmé, z ktorého zdroja údaj pochádza a za akých podmienok bol získaný. Evidencia provenancie preto neslúži iba na technickú auditovateľnosť, ale aj na transparentné odlíšenie údajov prevzatých z otvorene použiteľných služieb od údajov pochádzajúcich z licenčne obmedzeného prostredia. V tomto zmysle právne a etické požiadavky priamo podporujú návrhový princíp, podľa ktorého aplikácia nemá nekriticky zlúčiť všetky dostupné údaje do jednej neodlišenej vrstvy, ale má zachovať ich pôvod, kontext a režim použitia \cite{CrossrefMetadataRetrieval}.
